
\subsection*{3. Fibrés projectifs et fibrés en drapeaux}

\textbf{Proposition 3.1.} \uline{Soient} $S$ \uline{un schéma noethérien}, $\underline{E}$ \uline{un} $\underline{O}_S$\uline{-Module localement libre de rang} $n$, $\pi=(p_1,\ldots,p_k)$ \uline{une suite d'entiers tels que} $\sum_i p_i=n$, \uline{et} $X=\mathbb D_\pi(\underline{E})$ \uline{le schéma des drapeaux de type} $\pi$ \uline{de} $\underline{E}$. \uline{Alors l'homomorphisme canonique} :
\[
K_\bullet(S)\otimes_{K^{\bullet}(S)}K^{\bullet}(X)\longrightarrow K_\bullet(X)
\]
\uline{est un isomorphisme}.

On reprend les notations de VI 4.5. On note donc $\underline{L}_i$ les quotients canoniques de l'image inverse de $\underline{E}$ sur $X$, $\underline{F}_i$ le noyau de $\underline{L}_i\to \underline{L}_{i-1}$ ($i=1,\ldots,k$), et $\lambda_i^{(j)}$ la classe dans $K^{\bullet}(X)$ de $\wedge^i(\underline{F}_j)$. Alors les $\lambda_i^{(j)}$ engendrent $K^{\bullet}(X)$ comme $K^{\bullet}(S)$-algèbre (VI 4.6). De même, les images des $\lambda_i^{(j)}$ dans $K^{\bullet}(X_s)$ engendrent $K^{\bullet}(X_s)$ comme algèbre sur $K^{\bullet}(k(s))=\mathbb Z$ pour tout point $s$ de $S$. Comme $X_s$ est un schéma de drapeaux sur un corps, c'est un schéma régulier, donc l'homomorphisme $K^{\bullet}(X_s)\to K_\bullet(X_s)$ est un isomorphisme, si bien que l'homomorphisme $K^{\bullet}(X)\to K_\bullet(X_s)$ est surjectif. D'après 1.3, cela entraîne que l'homomorphisme:
\[
K_\bullet(S)\otimes_{K^{\bullet}(S)}K^{\bullet}(S)\longrightarrow K_\bullet(X)
\]
est surjectif.

Pour montrer l'injectivité, on peut se ramener d'abord au cas du fibré en drapeaux complet $\mathbb D(E)$. En effet, supposons que
\[
K_\bullet(S)\otimes_{K^{\bullet}(S)}K^{\bullet}(\mathbb D(\underline{E}))\longrightarrow K_\bullet(\mathbb D(\underline{E}))
\]
soit injectif. Comme $K_\bullet(S)\to K_\bullet(\mathbb D(\underline{E}))$ se factorise par $K_\bullet(X)$, cet
